% Part: first-order-logic
% Chapter: natural-deduction
% Section: proof-theoretic-notions

\documentclass[../../../include/open-logic-section]{subfiles}

\begin{document}

\iftag{FOL}
      {\olfileid{fol}{ntd}{ptn}}
      {\olfileid{pl}{ntd}{ptn}}

\olsection{સાબિતી-સૈદ્ધાંતિક ખ્યાલો}

\begin{editorial}
આ વિભાગ સ્વાભાવિક નિગમન માટે સાબિત કરી શકાય તેવા હોવાના સંબંધ અને
સુસંગતતાની વ્યાખ્યાઓ એકત્ર કરે છે.
\end{editorial}

\begin{explain}
આપણે અનેક મહત્વના અર્થાનુસારી ખ્યાલો (પ્રામાણ્ય, ફલિતતા અને સંતોષ્યતા)
વ્યાખ્યાયિત કર્યા છે; એ જ રીતે હવે તેમને અનુરૂપ \emph{સાબિતી-સૈદ્ધાંતિક
ખ્યાલો} વ્યાખ્યાયિત કરીએ છીએ. આ ખ્યાલો !!{structure}sમાં !!{sentence}sના
સંતોષનો આશરો લઈને નહિ, પરંતુ કેટલાક !!{sentence}sની બીજાં વાક્યોમાંથી
!!{derivability} અથવા !!{nonderivability}નો આશરો લઈને વ્યાખ્યાયિત થાય છે.
આ બંને પ્રકારના ખ્યાલો મળી જાય છે એ એક મહત્વની શોધ હતી. તેઓ મળે છે તે
જ \emph{યથાર્થતા} અને \emph{પૂર્ણતાના પ્રમેયો}નો વિષય છે.
\end{explain}


\begin{defn}[પ્રમેયો]
જો સ્વાભાવિક નિગમનમાં~$!A$ની એવી !!a{derivation} હોય જેમાં બધી ધારણાઓ
!!{discharged} હોય, તો !!{sentence}~$!A$ \emph{પ્રમેય} છે. $!A$
પ્રમેય હોય તો $\Proves !A$ અને ન હોય તો $\Proves/ !A$ લખીએ છીએ.
\end{defn}

\begin{defn}[!!^{derivability}]
!!^a{sentence} $!A$ એ !!{sentence}sના ગણ~$\Gamma$માંથી
\emph{!!{derivable}} છે, $\Gamma \Proves !A$, જો ફલિતવાક્ય~$!A$ ધરાવતી
એવી !!{derivation} હોય જેમાં દરેક ધારણા કાં તો !!{discharged} હોય અથવા
~$\Gamma$માં હોય. $!A$ એ $\Gamma$માંથી !!{derivable} ન હોય તો
$\Gamma \Proves/ !A$ લખીએ છીએ.
\end{defn}

\begin{defn}[સુસંગતતા]
!!{sentence}sનો ગણ~$\Gamma$ \emph{અસુસંગત} છે જો અને માત્ર જો
$\Gamma \Proves \lfalse$. જો $\Gamma$ અસુસંગત ન હોય, એટલે કે
$\Gamma \Proves/ \lfalse$, તો તેને \emph{સુસંગત} કહીએ છીએ.
\end{defn}

\begin{prop}[સ્વવાચકતા]
\ollabel{prop:reflexivity}
જો $!A \in \Gamma$ હોય, તો $\Gamma \Proves !A$.
\end{prop}

\begin{proof}
એકલી ધારણા $!A$ એ~$!A$ની એવી !!a{derivation} છે જેમાં દરેક
!!{undischarged} ધારણા (એટલે કે $!A$) ~ $\Gamma$માં છે.
\end{proof}


\begin{prop}[એકદિશવર્ધિતા]
\ollabel{prop:monotonicity}
જો $\Gamma \subseteq \Delta$ અને $\Gamma \Proves !A$ હોય, તો
$\Delta \Proves !A$.
\end{prop}

\begin{proof}
$\Gamma$માંથી $!A$ની કોઈ પણ !!{derivation} એ~$\Delta$માંથી $!A$ની પણ
!!a{derivation} છે.
\end{proof}

\begin{prop}[પરંપરિતતા]
\ollabel{prop:transitivity}
જો $\Gamma \Proves !A$ અને $\{!A\} \cup \Delta \Proves !B$ હોય, તો
$\Gamma \cup \Delta \Proves !B$.
\end{prop}

\begin{proof}
જો $\Gamma \Proves !A$ હોય, તો એવી !!a{derivation}~$\delta_0$ હોય છે
જેમાં~$!A$નું ફલિતવાક્ય અને બધી !!{undischarged} ધારણાઓ~$\Gamma$માં છે.
જો $\{!A\} \cup \Delta \Proves !B$ હોય, તો એવી
!!a{derivation}~$\delta_1$ હોય છે જેમાં~$!B$નું ફલિતવાક્ય અને બધી
!!{undischarged} ધારણાઓ~$\{!A\} \cup \Delta$માં છે. હવે જુઓ:
\begin{prooftree}
  \AxiomC{$\Delta, \Discharge{!A}{1}$}
  \RightLabel{$\delta_1$}
  \DeduceC{$!B$}
  \DischargeRule{\Intro{\lif}}{1}
  \UnaryInfC{$!A \lif !B$}
  \AxiomC{$\Gamma$}
  \RightLabel{$\delta_0$}
  \DeduceC{$!A$}
  \RightLabel{\Elim{\lif}}
  \BinaryInfC{$!B$}
\end{prooftree}
હવે બધી !!{undischarged} ધારણાઓ $\Gamma \cup \Delta$માં છે, તેથી આ
$\Gamma \cup \Delta \Proves !B$ બતાવે છે.
\end{proof}

$\Gamma = \{!A_1, !A_2, \ldots, !A_k\}$ સાન્ત ગણ હોય ત્યારે
$\Gamma \Proves !B$ માટે આપણે સરળ સંકેત
$!A_1,!A_2,\ldots,!A_k \Proves !B$ વાપરી શકીએ છીએ. ખાસ કરીને,
$!A \Proves !B$નો અર્થ $\{!A\} \Proves !B$ થાય છે.

નોંધો કે જો $\Gamma \Proves !A$ અને $!A \Proves !B$ હોય, તો
$\Gamma \Proves !B$. વળી, જો $!A_1, \dots, !A_n \Proves !B$ અને
દરેક~$i$ માટે $\Gamma \Proves !A_i$ હોય, તો $\Gamma \Proves !B$
એવું પણ ફલિત થાય છે.

\begin{prop}
\ollabel{prop:incons}
નીચેની શરતો સમકક્ષ છે.
    \begin{enumerate}
      \item \( \Gamma \) અસુસંગત છે.
      \item દરેક !!{sentence}~\( {!A} \) માટે \( \Gamma \Proves {!A} \).
      \item કોઈ !!{sentence}~\( {!A} \) માટે \( \Gamma \Proves {!A} \) અને \( \Gamma \Proves \lnot {!A} \).
    \end{enumerate}
\end{prop}

\begin{proof}
અભ્યાસપ્રશ્ન.
\end{proof}

\tagprob{FOL}
\begin{prob}
\olref[fol][ntd][ptn]{prop:incons} સાબિત કરો.
\end{prob}
\tagendprob

\tagprob{notFOL}
\begin{prob}
\olref[pl][ntd][ptn]{prop:incons} સાબિત કરો.
\end{prob}
\tagendprob

\begin{prop}[સઘનતા]
  \ollabel{prop:proves-compact}
  \begin{enumerate}
  \item જો $\Gamma \Proves !A$ હોય, તો એવો સાન્ત ઉપગણ
    $\Gamma_0 \subseteq \Gamma$ છે કે $\Gamma_0 \Proves !A$.
  \item જો $\Gamma$નો દરેક સાન્ત ઉપગણ સુસંગત હોય, તો $\Gamma$ સુસંગત છે.
  \end{enumerate}
\end{prop}

\begin{proof}
  \begin{enumerate}
    \item જો $\Gamma \Proves !A$ હોય, તો $\Gamma$માંથી~$!A$ની કોઈ
      !!a{derivation}~$\delta$ છે. $\delta$ની !!{undischarged}
      ધારણાઓના ગણને $\Gamma_0$ કહો. દરેક !!{derivation} સાન્ત હોવાથી
      $\Gamma_0$માં માત્ર સાન્ત સંખ્યામાં !!{sentence}s હોઈ શકે. તેથી
      $\delta$ એ સાન્ત $\Gamma_0 \subseteq \Gamma$માંથી~$!A$ની
      !!a{derivation} છે.
    \item $!A \ident \lfalse$ હોય એ વિશેષ કિસ્સામાં આ (1)નું
      પ્રતિવિધાન છે.
  \end{enumerate}
\end{proof}

\end{document}
